\section{Lecture 8: Logistic Regression and Classification}

\subsection{Motivation}

In the previous lecture, we studied linear regression, where outputs are continuous. We now turn to \emph{classification}, where outputs are discrete.

\medskip

\noindent
\textbf{Goal.}  
Given input $x \in \mathcal{X}$, predict a label:
\[
y \in \{0,1\}.
\]

\medskip

\noindent
A naive approach would be to apply linear regression and threshold the output. However, this lacks a probabilistic interpretation and can produce invalid predictions.

\medskip

\noindent
\textbf{Guiding idea.}  
Model the probability of class membership directly.

\subsection{Probabilistic Interpretation}

We model the conditional probability:
\[
p(y = 1 \mid x) = \sigma(w^\top x),
\]
where $\sigma$ is the sigmoid function:
\[
\sigma(z) = \frac{1}{1 + e^{-z}}.
\]

\medskip

\noindent
\textbf{Properties.}
\begin{itemize}
    \item $\sigma(z) \in (0,1)$
    \item Monotonically increasing
    \item Maps linear scores to probabilities
\end{itemize}

\medskip

\noindent
\textbf{Odds interpretation.}
\[
\frac{p(y=1 \mid x)}{p(y=0 \mid x)} = e^{w^\top x}.
\]

\medskip

\noindent
Taking logarithms:
\[
\log \frac{p(y=1 \mid x)}{p(y=0 \mid x)} = w^\top x.
\]

\medskip

\noindent
\textbf{Insight.}  
Logistic regression models the log-odds as a linear function of the input.

\subsection{Cross-Entropy Loss}

Given predictions $p_i = \sigma(w^\top x_i)$, we define the loss:
\[
\ell(w) = - \sum_{i=1}^n \left[
y_i \log p_i + (1 - y_i)\log(1 - p_i)
\right].
\]

\medskip

\noindent
\textbf{Interpretation.}
\begin{itemize}
    \item Penalizes confident incorrect predictions
    \item Encourages correct probability estimates
\end{itemize}

\medskip

\noindent
\textbf{Gradient.}
\[
\nabla_w \ell(w) = \sum_{i=1}^n (p_i - y_i)x_i.
\]

\medskip

\noindent
\textbf{Properties.}
\begin{itemize}
    \item Convex objective
    \item No closed-form solution
    \item Solved using iterative methods
\end{itemize}

\medskip

\noindent
\textbf{Insight.}  
Cross-entropy aligns optimization with probabilistic modeling.

\subsection{Maximum Likelihood Perspective}

We assume:
\[
y_i \sim \text{Bernoulli}(p_i), \quad p_i = \sigma(w^\top x_i).
\]

\medskip

\noindent
\textbf{Likelihood.}
\[
p(y_i \mid x_i, w) = p_i^{y_i}(1 - p_i)^{1 - y_i}.
\]

\medskip

\noindent
\textbf{Log-likelihood.}
\[
\sum_{i=1}^n \left[
y_i \log p_i + (1 - y_i)\log(1 - p_i)
\right].
\]

\medskip

\noindent
\textbf{Result.}  
Maximizing likelihood is equivalent to minimizing cross-entropy loss.

\medskip

\noindent
\textbf{Interpretation.}
\begin{itemize}
    \item Learning corresponds to estimating model parameters under a probabilistic model
\end{itemize}

\medskip

\noindent
\textbf{Insight.}  
Logistic regression arises naturally from maximum likelihood estimation.

\subsection{Decision Boundaries}

The predicted class is:
\[
\hat{y} =
\begin{cases}
1 & \text{if } p(y=1 \mid x) \geq 0.5 \\
0 & \text{otherwise}
\end{cases}
\]

\medskip

\noindent
Since $\sigma(z) \geq 0.5$ when $z \geq 0$, this becomes:
\[
w^\top x = 0.
\]

\medskip

\noindent
\textbf{Decision boundary.}
\begin{itemize}
    \item A hyperplane in $\mathbb{R}^d$
\end{itemize}

\medskip

\noindent
\textbf{Geometry.}
\begin{itemize}
    \item Points are classified based on which side of the hyperplane they lie
\end{itemize}

\medskip

\noindent
\textbf{Extension.}
\begin{itemize}
    \item Nonlinear boundaries can be obtained using feature maps $\phi(x)$
\end{itemize}

\medskip

\noindent
\textbf{Insight.}  
Logistic regression produces linear decision boundaries in feature space.

\subsection{A Unifying Perspective}

Logistic regression integrates key ideas from previous lectures:

\begin{itemize}
    \item \textbf{Function approximation:} linear model in feature space
    \item \textbf{Probability:} modeling conditional distributions
    \item \textbf{Optimization:} minimizing cross-entropy
    \item \textbf{Geometry:} linear decision boundaries
\end{itemize}

\medskip

\noindent
\textbf{Core idea.}  
Classification can be framed as probabilistic modeling combined with optimization.

\medskip

\noindent
\textbf{Key insight.}  
Logistic regression is a bridge between linear models and more complex classification methods.

\subsection{Outlook}

In this lecture, we studied:

\begin{itemize}
    \item probabilistic modeling of classification,
    \item cross-entropy loss,
    \item maximum likelihood estimation,
    \item geometric interpretation via decision boundaries.
\end{itemize}

\medskip

\noindent
However, linear models are limited in their expressivity.

\medskip

\noindent
In the next lecture, we introduce feature maps and kernel methods, which enable nonlinear modeling while retaining linear structure.

\medskip

\noindent
\textbf{Guiding principle.}  
By combining representations with linear models, we can capture complex patterns in data.